\FloatBarrier
\section{Orbital States in the Effective-Mass Approximation}
\label{app:effective-mass}
\label{sec:confinement-single-electron}

Section~\ref{sec:quantum-dot-introduction} uses
Eq.~\eqref{eq:dot-orbital-eigenproblem} to describe a localized orbital.
This appendix derives that equation from the crystal problem and then
shows how to obtain bound orbitals from a specified confining potential.


Our goal is to derive a Schr\"odinger equation for a confined
conduction electron without resolving its wave function inside
every crystal cell. The lattice will enter through an effective
mass, while the dot potential will determine a slowly varying
envelope. The main issue is how to derive that envelope equation
when the envelope itself is unknown.

We first expand the microscopic state in a complete basis with
ordinary numerical coefficients. Projection gives equations for
those coefficients; a Fourier transform then gives coupled
equations for the envelopes. Finally, we express the small
other-band components in terms of the conduction-band envelope
and substitute them back. This is a fixed-reference
Luttinger--Kohn construction~\cite{LuttingerKohn1955}.
We work in one dimension to expose each step and give the
three-dimensional result at the end.

\subsection{The crystal reference and the confined state}
\label{sec:em-crystal-reference}

Consider the spinless, one-electron Hamiltonian
\begin{equation}
  H_{\mathrm{free}}=\frac{p^2}{2m_e}+V_{\mathrm{lat}}(x),
  \qquad p=-i\hbar\partial_x,
  \qquad V_{\mathrm{lat}}(x+a)=V_{\mathrm{lat}}(x),
  \label{eq:em-crystal-hamiltonian}
\end{equation}
where $a$ is the lattice period, $m_e$ is the vacuum electron
mass, and $V_{\mathrm{lat}}$ is a real local potential energy.
Here ``free'' means without dot
confinement, not without the crystal lattice. We take this
crystal model and its Bloch eigenstates as given:
\begin{equation}
  H_{\mathrm{free}}\chi_{n,k}=E_n(k)\chi_{n,k},
  \qquad
  \chi_{n,k}(x)=e^{ikx}u_{n,k}(x),
  \qquad u_{n,k}(x+a)=u_{n,k}(x).
  \label{eq:em-bloch-states}
\end{equation}
Throughout this appendix, $n$ and $m$ label crystal bands, not
the discrete orbitals of a quantum dot.

Choose a nondegenerate conduction-band minimum at $k_0$, with
band index $\mathrm c$. We abbreviate
$E_n=E_n(k_0)$ and $\chi_n=\chi_{n,k_0}$.
Adding a dot potential gives the eigenvalue problem
\begin{equation}
  H=H_{\mathrm{free}}+V_{\mathrm{conf}}(x),
  \qquad H\psi=E\psi.
  \label{eq:em-confined-problem}
\end{equation}
If the potential supports a bound state, $\psi$ is localized,
whereas the reference state $\chi_{\mathrm c}$ is extended.
The eventual approximation concerns small changes in band
character near $k_0$; it does not require the localized state
to be close to the extended state in norm.

Locally, wherever $\chi_{\mathrm c}\ne0$, one can write
$\psi=\chi_{\mathrm c}F$ exactly. Product differentiation gives
\begin{equation}
  \left[-\frac{\hbar^2}{2m_e}\partial_x^2
        -\frac{\hbar^2}{m_e}\frac{\chi_{\mathrm c}'}{\chi_{\mathrm c}}
             \partial_x+V_{\mathrm{conf}}\right]F
       =(E-E_{\mathrm c})F.
  \label{eq:em-exact-quotient}
\end{equation}
Here a prime denotes an $x$ derivative. This is not yet an
effective-mass equation: $F$ can contain microscopic variations,
and the quotient can be singular at nodes of $\chi_{\mathrm c}$.
Dropping the coefficient involving $\chi_{\mathrm c}'$ would
leave $m_e$, not produce a material-dependent mass. We therefore
resolve the microscopic components before making an approximation.

\subsection{A complete basis at a fixed reference wave vector}
\label{sec:em-complete-basis}

\paragraph{Completeness within one cell.}
Substituting the Bloch form into Eq.~\eqref{eq:em-bloch-states}
gives an eigenproblem for the periodic part:
\begin{equation}
  H(k)u_{n,k}=E_n(k)u_{n,k},
  \qquad H(k)=\frac{(p+\hbar k)^2}{2m_e}+V_{\mathrm{lat}}.
  \label{eq:em-cell-hamiltonian}
\end{equation}
At fixed $k_0$, this is one operator on $[0,a]$ with periodic
boundary conditions on both the function and its derivative.
Fixing $k_0$ fixes an operator parameter; it does not select
one of its eigenstates. The band index enumerates all of them.

For a sufficiently regular real periodic potential, this
self-adjoint cell operator has a discrete spectrum and a complete
orthonormal eigenbasis. The mathematical reason is that it has
compact resolvent: for a constant $C$ chosen to make $H(k_0)+C$
positive, the inverse $(H(k_0)+C)^{-1}$ is compact. In the
zero-potential Fourier basis its eigenvalues tend to zero as
the Fourier wave vector grows; a bounded periodic potential
preserves this property. The spectral theorem then supplies
the complete basis, rather than self-adjointness alone.
We use the cell-average normalization
\begin{equation}
  \langle u_{n,k_0},u_{m,k_0}\rangle_{\mathrm{cell}}
    \equiv\frac1a\int_0^a u_{n,k_0}^*(x)u_{m,k_0}(x)\,dx
    =\delta_{nm}.
  \label{eq:em-cell-normalization}
\end{equation}
For example, when $V_{\mathrm{lat}}=0$, the periodic eigenfunctions
can be chosen as $e^{iGx}$, with $G=2\pi\ell/a$ and integer
$\ell$. These are precisely the Fourier basis of a cell.
The periodic potential mixes these modes; its complete eigenbasis
still spans the same function space.

The periodic parts at any other wave vector belong to this
same space. Consequently,
\begin{equation}
  u_{j,k_0+q}(x)=\sum_n b_{nj}(q)u_{n,k_0}(x),
  \qquad
  b_{nj}(q)=\langle u_{n,k_0},u_{j,k_0+q}\rangle_{\mathrm{cell}}.
  \label{eq:em-cell-basis-change}
\end{equation}
Here $j$ is another band index. With all bands retained, this
is an exact change between two orthonormal cell bases.
It does not claim that $u_{j,k_0+q}$ is an eigenfunction of $H(k_0)$.

\paragraph{Completeness across the crystal.}
The fixed-$k_0$ states $\chi_n$ alone do not span arbitrary
wave functions over the entire crystal. To include changes from
cell to cell, introduce
\begin{equation}
  B_{nq}(x)=\frac{1}{\sqrt{2\pi}}e^{iqx}\chi_n(x)
    =\frac{1}{\sqrt{2\pi}}e^{i(k_0+q)x}u_{n,k_0}(x),
  \qquad q\in\mathcal Q=[-\pi/a,\pi/a).
  \label{eq:em-global-basis}
\end{equation}
The interval $k_0+\mathcal Q$ contains one representative of
every crystal quasi-momentum modulo $2\pi/a$.
Each $B_{nq}$ has quasi-momentum $k_0+q$, but is generally not
an energy eigenstate there. Equation~\eqref{eq:em-cell-basis-change}
expresses every true Bloch state as a combination of these
functions at the same $q$, establishing completeness.

To verify orthogonality, expand the periodic product as
\begin{equation}
  u_{n,k_0}^*u_{m,k_0}=\sum_G S_{nm}(G)e^{iGx},
  \qquad
  S_{nm}(G)=\frac1a\int_0^a
       u_{n,k_0}^*u_{m,k_0}e^{-iGx}\,dx.
  \label{eq:em-cell-product}
\end{equation}
Integration over the whole line then gives
\begin{equation}
  \int_{\mathbb R}B_{nq}^*(x)B_{mq'}(x)\,dx
    =\sum_G S_{nm}(G)\delta(q'-q+G)
    =\delta_{nm}\delta(q-q').
  \label{eq:em-global-orthogonality}
\end{equation}
The last equality follows because two points of the half-open
interval $\mathcal Q$ cannot differ by a nonzero reciprocal
lattice vector, and $S_{nm}(0)=\delta_{nm}$.
The continuous basis is delta-normalized; a finite periodic
crystal gives the corresponding discrete sums.

\subsection{Unknown envelopes and an exact projection}
\label{sec:em-exact-projection}

We can now expand the unknown state without treating a
position-dependent function as a numerical coefficient:
\begin{equation}
  \psi(x)=\sum_n\int_{\mathcal Q}dq\,A_n(q)B_{nq}(x),
  \qquad
  A_n(q)=\int_{\mathbb R}B_{nq}^*(x)\psi(x)\,dx.
  \label{eq:em-state-expansion}
\end{equation}
For fixed labels $n,q$, $A_n(q)$ is an unknown complex number
independent of $x$. Define its Fourier transform by
\begin{equation}
  f_n(x)=\frac{1}{\sqrt{2\pi}}\int_{\mathcal Q}dq\,
                          A_n(q)e^{iqx},
  \qquad \phi(x)\equiv f_{\mathrm c}(x).
  \label{eq:em-envelope-definition}
\end{equation}
Rearranging the exact expansion gives
\begin{equation}
  \psi=\sum_n\chi_nf_n
       =\chi_{\mathrm c}\phi+\sum_{n\ne\mathrm c}\chi_nf_n.
  \label{eq:em-multiband-reconstruction}
\end{equation}
Both $\phi$ and the other $f_n$ vary with position and are unknown.
Their Fourier restriction to the same interval $\mathcal Q$
is essential: arbitrary functions multiplying the $\chi_n$
would give a nonunique decomposition. Given $\psi$, the complete
basis instead determines every $A_n$, and hence every $f_n$,
uniquely. No assumption of slow variation has yet been made;
that will require concentration near $q=0$.

Projection must use Eq.~\eqref{eq:em-state-expansion} and
whole-crystal inner products. At a fixed $x$, the values
$\chi_n(x)$ are just complex numbers; one cannot compare
their function-valued coefficients as independent coordinates.
Nor can $f_m(x)$ be pulled out of the cell integral
$\int_0^a\chi_n^*(x)\chi_m(x)f_m(x)\,dx$.

\paragraph{The crystal part of the projected equation.}
Define the momentum matrix elements at the reference wave vector:
\begin{equation}
  p_{nm}=\frac1a\int_0^a
       u_{n,k_0}^*(p+\hbar k_0)u_{m,k_0}\,dx
       =\frac1a\int_0^a\chi_n^*p\chi_m\,dx,
  \qquad p_{mn}=p_{nm}^*.
  \label{eq:em-momentum-elements}
\end{equation}
These are numbers, with units of momentum.
Acting on a basis function and differentiating the factor
$e^{iqx}$ gives
\begin{equation}
  H_{\mathrm{free}}B_{mq}
    =\left(E_m+\frac{\hbar^2q^2}{2m_e}\right)B_{mq}
       +\frac{\hbar q}{m_e}\sum_n p_{nm}B_{nq}.
  \label{eq:em-crystal-action}
\end{equation}
The sum arises by expanding $(p+\hbar k_0)u_{m,k_0}$ in the
complete cell basis. There are no unknown envelopes in that
cell expansion.

\paragraph{The confining potential and the full equation.}
Let
\begin{equation}
  V_{nm}(q,q')=\int_{\mathbb R}B_{nq}^*(x)
                           V_{\mathrm{conf}}(x)B_{mq'}(x)\,dx.
  \label{eq:em-potential-kernel}
\end{equation}
Multiplying $H\psi=E\psi$ by $B_{nq}^*$ and integrating now
produces the exact coupled equations
\begin{equation}
  \begin{aligned}
    0={}&\left(E_n+\frac{\hbar^2q^2}{2m_e}-E\right)A_n(q)
       +\frac{\hbar q}{m_e}\sum_m p_{nm}A_m(q)\\
       &+\sum_m\int_{\mathcal Q}dq'\,V_{nm}(q,q')A_m(q').
  \end{aligned}
  \label{eq:em-exact-amplitude-equations}
\end{equation}
This is a representation of the microscopic problem, not yet
a reduced model. In particular, both changes of quasi-momentum
and mixing of the reference-band components are present.

\subsection{The smooth-potential approximation}
\label{sec:em-smooth-potential}

We now select low-energy states concentrated near $q=0$ and
a potential varying on lengths much greater than $a$.
The reason a smooth potential becomes approximately diagonal
in the reference-band index can be seen from its Fourier transform.
With the convention
\begin{equation}
  \widetilde V(s)=\frac1{2\pi}\int_{\mathbb R}dx\,
                  e^{-isx}V_{\mathrm{conf}}(x),
  \label{eq:em-potential-fourier}
\end{equation}
Eq.~\eqref{eq:em-cell-product} gives the exact kernel
\begin{equation}
  V_{nm}(q,q')=\sum_G S_{nm}(G)
                           \widetilde V(q-q'-G).
  \label{eq:em-potential-reciprocal-sum}
\end{equation}
For small $q,q'$, every $G\ne0$ samples a Fourier component
at a lattice-scale wave vector. Smooth confinement has little
weight there. Retaining $G=0$ gives
\begin{equation}
  V_{nm}(q,q')\simeq\delta_{nm}\widetilde V(q-q').
  \label{eq:em-smooth-kernel}
\end{equation}
Thus it is the potential's lack of large-wave-vector components,
not the constancy of an unknown envelope, that justifies this step.

To obtain a local differential equation, we also neglect weight
near the ends of $\mathcal Q$ and extend the low-wave-vector
Fourier integrals to the real line. The convolution theorem then
turns the potential term into $V_{\mathrm{conf}}f_n$.
At finite $\mathcal Q$, it would instead be the Fourier projection
of that product back into $\mathcal Q$; replacing it by the
local product is valid when the discarded weight is negligible.
This extension is an approximation, not a second exact basis expansion.

The remaining Fourier identities are
\[
  \frac1{\sqrt{2\pi}}\int dq\,qA_n(q)e^{iqx}=-if_n'(x),
  \qquad
  \frac1{\sqrt{2\pi}}\int dq\,q^2A_n(q)e^{iqx}=-f_n''(x).
\]
Using them in Eq.~\eqref{eq:em-exact-amplitude-equations} gives
the leading long-wavelength envelope system
\begin{equation}
  (\Delta_n+V_{\mathrm{conf}}-\epsilon)f_n
     -\frac{\hbar^2}{2m_e}f_n''
     -\frac{i\hbar}{m_e}\sum_m p_{nm}f_m'=0,
  \quad
  \Delta_n=E_n-E_{\mathrm c},\quad \epsilon=E-E_{\mathrm c}.
  \label{eq:em-coupled-envelopes}
\end{equation}
Here $\Delta_{\mathrm c}=0$. All envelopes are still unknown
functions. Equation~\eqref{eq:em-coupled-envelopes} was obtained
by transforming projected equations for numerical amplitudes,
not by pulling functions out of a cell integral.

\subsection{Eliminating the other-band envelopes}
\label{sec:em-band-elimination}

The $n=\mathrm c$ component describes the desired low-energy
state. The other components have nonzero reference energy
separations $\Delta_n$. When those separations dominate the
envelope energy and derivative couplings, their amplitudes can
be expressed perturbatively in terms of $\phi=f_{\mathrm c}$.
They need not be known beforehand.

Let $\ell$ be a typical envelope length and
$\Delta_{\min}=\min_{n\ne\mathrm c}|\Delta_n|>0$.
We require $\ell\gg a$, envelope kinetic and local energy scales
small compared with $\Delta_{\min}$, and weak derivative coupling
to the excluded bands. In particular,
\begin{equation}
  \frac{\hbar^2}{2m_e\ell^2},\ |V_{\mathrm{conf}}-\epsilon|
        \ll\Delta_{\min},
  \qquad
  \frac{\hbar|p_{n\mathrm c}|}
       {m_e\ell|\Delta_n|}\ll1
       \quad(n\ne\mathrm c),
  \label{eq:em-elimination-conditions}
\end{equation}
with the analogous smallness of derivative couplings within
the excluded subspace. Local energy conditions concern the
region appreciably occupied by the state. A nearby degenerate
band must be retained, not eliminated with a small denominator.
The $\Delta_n$ can be positive or negative: ``other bands''
includes bands below the chosen conduction band.
We retain terms through second order in slow derivatives and
to leading order in the ratio of envelope energy to band separation.

At the minimum, the conduction-band diagonal momentum vanishes.
Differentiating the cell eigenvalue equation and using normalization
gives the Hellmann--Feynman identity
\begin{equation}
  \left.\frac{dE_{\mathrm c}(k)}{dk}\right|_{k_0}
      =\left\langle u_{\mathrm c,k_0},
        \left.\frac{\partial H(k)}{\partial k}\right|_{k_0}
          u_{\mathrm c,k_0}\right\rangle_{\mathrm{cell}}
      =\frac{\hbar}{m_e}p_{\mathrm{cc}}=0.
  \label{eq:em-zero-group-velocity}
\end{equation}
Consequently, the conduction component of
Eq.~\eqref{eq:em-coupled-envelopes} is
\begin{equation}
  -\frac{\hbar^2}{2m_e}\phi''+V_{\mathrm{conf}}\phi
    -\frac{i\hbar}{m_e}\sum_{n\ne\mathrm c}p_{\mathrm c n}f_n'
       =\epsilon\phi.
  \label{eq:em-conduction-envelope}
\end{equation}

For $n\ne\mathrm c$, the leading source in its equation is
the derivative of $\phi$. To first order in slow spatial
variation, balance that source against the band separation:
\begin{equation}
  \Delta_nf_n-\frac{i\hbar}{m_e}p_{n\mathrm c}\phi'\simeq0,
  \qquad
  f_n\simeq\frac{i\hbar p_{n\mathrm c}}{m_e\Delta_n}\phi'.
  \label{eq:em-eliminated-envelopes}
\end{equation}
Derivatives of these already small components give higher
orders in their own equations; their local energy corrections
are suppressed by Eq.~\eqref{eq:em-elimination-conditions}.
The displayed term is enough to determine the leading
second-derivative correction in the conduction equation.
Although $f_n$ is small, setting it to zero before this
substitution would discard the mass correction.

The crystal parameters are position independent here. Substituting
$f_n'$ from Eq.~\eqref{eq:em-eliminated-envelopes} into
Eq.~\eqref{eq:em-conduction-envelope}, with
$p_{\mathrm c n}=p_{n\mathrm c}^*$, yields
\begin{equation}
  \left[-\frac{\hbar^2}{2m_e}
       +\frac{\hbar^2}{m_e^2}\sum_{n\ne\mathrm c}
                 \frac{|p_{n\mathrm c}|^2}{\Delta_n}\right]\phi''
       +V_{\mathrm{conf}}\phi=\epsilon\phi.
  \label{eq:em-mass-before-definition}
\end{equation}
Define the effective mass by collecting the coefficient of
$-\hbar^2\phi''/2$:
\begin{equation}
  \boxed{\frac1{m^*}=\frac1{m_e}
      +\frac{2}{m_e^2}\sum_{n\ne\mathrm c}
                    \frac{|p_{n\mathrm c}|^2}{E_{\mathrm c}-E_n}}.
  \label{eq:em-mass-sum}
\end{equation}
The desired single-envelope equation is therefore
\begin{equation}
  \boxed{\left[-\frac{\hbar^2}{2m^*}\partial_x^2
           +V_{\mathrm{conf}}(x)\right]\phi(x)
          =\epsilon\phi(x)}.
  \label{eq:em-effective-equation}
\end{equation}
The other-band amplitudes describe components of the same
stationary state in the reference basis, not real transitions
that populate separate excited states in time. Their coupling
back to the conduction component changes its kinetic energy.

\subsection{Band curvature and reconstruction of the microscopic state}
\label{sec:em-curvature-reconstruction}

\paragraph{Why the mass equals the band curvature.}
Set $V_{\mathrm{conf}}=0$ and examine one small wave-vector
offset $q$. In the fixed cell basis,
\begin{equation}
  H(k_0+q)=H(k_0)+\frac{\hbar q}{m_e}(p+\hbar k_0)
                       +\frac{\hbar^2q^2}{2m_e}.
  \label{eq:em-kp-expansion}
\end{equation}
The first-order admixture of band $n\ne\mathrm c$ is
$\hbar q p_{n\mathrm c}/[m_e(E_{\mathrm c}-E_n)]$.
Feeding it back into the conduction component gives
\begin{equation}
  E_{\mathrm c}(k_0+q)-E_{\mathrm c}
     =\frac{\hbar^2q^2}{2m_e}
       +\frac{\hbar^2q^2}{m_e^2}
          \sum_{n\ne\mathrm c}
             \frac{|p_{n\mathrm c}|^2}{E_{\mathrm c}-E_n}
       +O(q^3).
  \label{eq:em-derived-dispersion}
\end{equation}
There is no linear term because $p_{\mathrm{cc}}=0$.
Comparison with Eq.~\eqref{eq:em-mass-sum} gives
\begin{equation}
  \frac1{m^*}=\frac1{\hbar^2}
       \left.\frac{d^2E_{\mathrm c}(k)}{dk^2}\right|_{k_0}.
  \label{eq:em-curvature-mass}
\end{equation}
Thus band-curvature matching and elimination of the reference-band
components give the same mass. Matching identifies the answer
from known dispersion; the projection derivation explains
which microscopic amplitudes generate it.

\paragraph{What the envelope approximates.}
After solving Eq.~\eqref{eq:em-effective-equation} with the
appropriate bound-state boundary conditions, reconstruct
\begin{equation}
  \psi(x)\simeq\chi_{\mathrm c}(x)\phi(x)
      +\sum_{n\ne\mathrm c}\chi_n(x)
          \frac{i\hbar p_{n\mathrm c}}{m_e\Delta_n}\phi'(x).
  \label{eq:em-first-corrector}
\end{equation}
The leading approximation is $\psi\simeq\chi_{\mathrm c}\phi$;
the displayed derivative term is its first microscopic correction.
Where division by $\chi_{\mathrm c}$ is permitted,
\begin{equation}
  F=\frac{\psi}{\chi_{\mathrm c}}
    \simeq\phi+\sum_{n\ne\mathrm c}
       \frac{\chi_n}{\chi_{\mathrm c}}
       \frac{i\hbar p_{n\mathrm c}}{m_e\Delta_n}\phi'.
  \label{eq:em-quotient-reconstruction}
\end{equation}
This gives the relation between the exact multiplicative
correction $F$ and the leading smooth envelope $\phi$.
They are not exactly the same unknown. Equation~\eqref{eq:em-first-corrector}
also remains meaningful at nodes where the quotient does not.

The complete-basis normalization satisfies
$\|\psi\|^2=\sum_n\int_{\mathbb R}|f_n(x)|^2\,dx$.
For weak admixture, $\int|\phi|^2dx=1$ is the leading normalization;
retaining the other components changes it at second order
in their relative amplitudes.

\subsection{Three dimensions and the range of validity}
\label{sec:em-three-dimensions}

In three dimensions, the linear coupling is
$\hbar\mathbf q\cdot\mathbf p_{nm}/m_e$. Define its Cartesian
components by
\[
  p_{nm}^{\alpha}=\left\langle u_{n,\mathbf k_0},
       (-i\hbar\partial_{r_\alpha}+\hbar k_{0,\alpha})
                u_{m,\mathbf k_0}\right\rangle_{\mathrm{cell}},
  \qquad \alpha=x,y,z,
\]
where the cell inner product is now the integral divided by
the cell volume. The same elimination gives a real symmetric
inverse-mass tensor:
\begin{equation}
  (\mathbf M^{-1})_{\alpha\beta}
    =\frac{\delta_{\alpha\beta}}{m_e}
      +\frac{2}{m_e^2}\operatorname{Re}
          \sum_{n\ne\mathrm c}
          \frac{p_{\mathrm c n}^{\alpha}p_{n\mathrm c}^{\beta}}
               {E_{\mathrm c}-E_n}
    =\frac1{\hbar^2}
       \left.\frac{\partial^2E_{\mathrm c}}
                    {\partial k_\alpha\partial k_\beta}
       \right|_{\mathbf k_0}.
  \label{eq:em-mass-tensor}
\end{equation}
Accordingly,
\begin{equation}
  \left[-\frac{\hbar^2}{2}\sum_{\alpha,\beta}
         (\mathbf M^{-1})_{\alpha\beta}
              \partial_{r_\alpha}\partial_{r_\beta}
        +V_{\mathrm{conf}}(\mathbf r)\right]\phi
      =\epsilon\phi.
  \label{eq:em-three-dimensional-equation}
\end{equation}
At a nondegenerate minimum with positive curvature, the tensor
is positive definite. This is the constant inverse-mass tensor used
in Eq.~\eqref{eq:dot-orbital-eigenproblem}. In rotated coordinates
$s_j$ along its principal axes, with eigenvalues $1/m_j$ for
$j=1,2,3$, the kinetic operator reduces to
$-\sum_{j=1}^{3}\hbar^2\partial_{s_j}^2/(2m_j)$.
These axes need not coincide with the device coordinates $x,y,z$.

The complete expansion and its projection were exact within
the stated local, spinless crystal model. The reduced equation
additionally assumes low wave-vector offsets, smooth potential
coupling, and weak admixture of separated bands, retaining the
leading quadratic dispersion. Strong confinement can require
higher dispersion terms; near-degenerate bands require coupled
envelopes rather than the scalar mass formula.

For silicon, the restriction to one minimum is a separate
single-valley assumption. Other valleys are not removed by an
energy denominator at that minimum. Atomic-scale interfaces can
couple valleys and require a multivalley description; a bulk
effective-mass tensor alone does not determine valley splitting~\cite{Zwanenburg2013}.
Spatially changing materials, nonlocal crystal operators, and
spin--orbit terms also require extensions of the model used here.

\subsection{From a three-dimensional well to bound orbitals}
\label{sec:device-bound-orbitals}

We first consider one valley $\nu$ and neglect coupling to other
valleys. The potential $V_{\mathrm{conf},\nu}$ is supplied by the
electrostatic calculation in Appendix~\ref{app:device-confinement};
we now find its bound states. Assuming smooth confinement on the lattice scale,
energies near this conduction-band minimum, and weak mixing with
other bands, the envelope Hamiltonian is
\begin{equation}
 h_\nu=-\frac{\hbar^2}{2}\nabla\cdot
                  [\mathbf M_\nu^{-1}(\mathbf r)\nabla]
              +V_{\mathrm{conf},\nu}(\mathbf r),
 \qquad
 h_\nu F_{n\nu}=\epsilon_{n\nu}F_{n\nu}.
 \label{eq:device-envelope-hamiltonian}
\end{equation}
$\mathbf M_\nu$ is the effective-mass tensor, and $F_{n\nu}$ is the
slowly varying spatial envelope of orbital $n$.
For a constant tensor this is the Hamiltonian of
Eq.~\eqref{eq:dot-orbital-eigenproblem}, whose relation to the microscopic
crystal problem has been derived above.
A spatially varying mass is written in divergence form so that
the kinetic operator respects the appropriate interface matching.

Strong vertical confinement can simplify this three-dimensional
problem. Suppose, for illustration, that one valley has a separable
potential $V_z(z)+V_\parallel(x,y)$. Assume additionally that the
mass tensor is constant and its principal axes align with the device
coordinates $x,y,z$, with corresponding masses $m_x,m_y,m_z$. First solve
\begin{equation}
 \left[-\frac{\hbar^2}{2m_z}\frac{d^2}{dz^2}+V_z(z)\right]\zeta_0(z)
       =\epsilon_z\zeta_0(z).
 \label{eq:device-vertical-state}
\end{equation}
If the vertical excitation gap is large compared with the energies
being used, the electron stays in this state. Writing
$F_n(x,y,z)=f_n(x,y)\zeta_0(z)$ then gives
\begin{equation}
 \left[-\frac{\hbar^2}{2m_x}\partial_x^2
       -\frac{\hbar^2}{2m_y}\partial_y^2
       +V_\parallel(x,y)\right]f_n
        =(\epsilon_n-\epsilon_z)f_n.
 \label{eq:device-lateral-state}
\end{equation}
The lateral potential governs motion in the material plane
within the chosen vertical state. A well in this potential
supports localized dot orbitals.

In a general device, a gate voltage changes both the vertical
state and lateral confinement. The full three-dimensional
equation includes this coupled response. A reduced-dimensional
model is useful when the vertical excitation gap is sufficiently
large and its dependence on the lateral coordinates is treated
consistently.

The resulting orbital energies and tails have different uses.
The energy gap to an excited orbital determines whether that orbital
can be omitted from a low-energy model. The tails extending into a
barrier determine coupling to another dot. Both must be resolved
numerically. Artificial hard walls are acceptable only where moving
them farther away no longer changes the states or couplings of interest.
